% GENERATED FILE. Edit canonical lessons, then run npm run generate:books.
% canonical-source-hash: fnv1a64:36e3a11e9c1c58bc
% canonical-lessons: UR-C311-harana, UR-C311-pachharna, UR-C311-fateh, UR-C311-shahid, UR-C311-sulah

\chapter{Victory and Defeat}
\label{ch:ur-victory-and-defeat}
\hlchaptermodality{311}

\begin{chapteropening}
\textbf{By the end of this chapter:} \emph{I can say who won a battle and how it ended.}

Complete the last lesson of chapter 311: I can say who won a battle and how it ended.
\end{chapteropening}

\section[\texorpdfstring{\ur{ہرانا} (harānā)}{ہرانا (harānā)}]{\ur{ہرانا} (harānā) — to defeat}

\label{lesson:UR-C311-harana}



\begin{quote}
Before the new one: say the Urdu for a bow for arrows, then the Urdu for an army on the march.
\end{quote}

\subsection*{You'll want to know: \ur{ہرانا}}
\textbf{\ur{ہرانا}} — \emph{harānā} — \textquotedblleft{}to defeat\textquotedblright{}.

\begin{quote}
\ur{ہماری} \ur{فوج} \ur{نے} \ur{انہیں} \ur{ہرایا۔}
\end{quote}

\begin{quote}
\emph{hamārī fauj ne unheṅ harāyā.}
\end{quote}

\begin{quote}
Our army defeated them.
\end{quote}

Past, with \ur{نے} (he or she …): \ur{اس} \ur{نے} \ur{بادشاہ} \ur{کو} \ur{ہرایا}, \emph{us ne bādshāh ko harāyā}.

\begin{cousinweb}
\textbf{\ur{ہرانا}} (\emph{harānā}) is inherited: CDIAL lists the Hindi cousins \emph{hārnā}, \textquotedblleft{}to lose\textquotedblright{}, and \emph{harānā}, \textquotedblleft{}to defeat\textquotedblright{}, under Sanskrit \emph{hārayati} (CDIAL 14061). It is the causing partner of \textbf{\ur{ہارنا}}.
\end{cousinweb}

\begin{grammarlens}[title={What you've built}]
One more word for telling what happened.
\end{grammarlens}

\subsection*{Your turn}
\begin{itemize}
\raggedright
  \item \emph{Say it:} \emph{harānā}
  \item \emph{Say it:} \emph{harānā}, once more
  \item \emph{Say it:} \emph{hamārī fauj ne unheṅ harāyā.} — Our army defeated them.
  \item \emph{Recall:} say the Urdu for a bow for arrows, then the Urdu for an army on the march, then say \emph{harānā} again
\end{itemize}

\begin{tcolorbox}[breakable,colback=teal!4,colframe=teal!35!black,title={Before you move on}]
What does \ur{ہرانا} mean? (\textquotedblleft{}To defeat\textquotedblright{}.) Say it once more.
\end{tcolorbox}
\section[\texorpdfstring{\ur{پچھاڑنا} (pachhāṛnā)}{پچھاڑنا (pachhāṛnā)}]{\ur{پچھاڑنا} (pachhāṛnā) — to throw down; to overpower}

\label{lesson:UR-C311-pachharna}



\begin{quote}
Before the new one: say the Urdu for an army on the march, then the Urdu for to defeat.
\end{quote}

\subsection*{You'll want to know: \ur{پچھاڑنا}}
\textbf{\ur{پچھاڑنا}} — \emph{pachhāṛnā} — \textquotedblleft{}to throw down; to overpower\textquotedblright{}.

\begin{quote}
\ur{اس} \ur{نے} \ur{سپاہی} \ur{کو} \ur{زمین} \ur{پر} \ur{پچھاڑا۔}
\end{quote}

\begin{quote}
\emph{us ne sipāhī ko zamīn par pachhāṛā.}
\end{quote}

\begin{quote}
He threw the soldier to the ground.
\end{quote}

Past, with \ur{نے} (he or she …): \ur{اس} \ur{نے} \ur{اسے} \ur{پچھاڑا}, \emph{us ne use pachhāṛā}.

\begin{cousinweb}
\textbf{\ur{پچھاڑنا}} (\emph{pachhāṛnā}) is inherited: CDIAL lists the Hindi cousin \emph{pachāṛnā}, \textquotedblleft{}to knock down\textquotedblright{}, under a form it reconstructs (CDIAL 8493).
\end{cousinweb}

\begin{grammarlens}[title={What you've built}]
One more word for telling what happened.
\end{grammarlens}

\subsection*{Your turn}
\begin{itemize}
\raggedright
  \item \emph{Say it:} \emph{pachhāṛnā}
  \item \emph{Say it:} \emph{pachhāṛnā}, once more
  \item \emph{Say it:} \emph{us ne sipāhī ko zamīn par pachhāṛā.} — He threw the soldier to the ground.
  \item \emph{Recall:} say the Urdu for an army on the march, then the Urdu for to defeat, then say \emph{pachhāṛnā} again
\end{itemize}

\begin{tcolorbox}[breakable,colback=teal!4,colframe=teal!35!black,title={Before you move on}]
What does \ur{پچھاڑنا} mean? (\textquotedblleft{}To throw down; to overpower\textquotedblright{}.) Say it once more.
\end{tcolorbox}
\section[\texorpdfstring{\ur{فتح} (fateh)}{فتح (fateh)}]{\ur{فتح} (fateh) — a victory}

\label{lesson:UR-C311-fateh}



\begin{quote}
Before the new one: say the Urdu for to defeat, then the Urdu for to throw down.
\end{quote}

\subsection*{You'll want to know: \ur{فتح}}
\textbf{\ur{فتح}} — \emph{fateh} — \textquotedblleft{}a victory\textquotedblright{}.

\begin{quote}
\ur{فتح} \ur{کے} \ur{بعد} \ur{جشن} \ur{ہوا۔}
\end{quote}

\begin{quote}
\emph{fateh ke bād jashn huā.}
\end{quote}

\begin{quote}
After the victory there was a celebration.
\end{quote}

Feminine.

\begin{cousinweb}
\textbf{\ur{فتح}} (\emph{fateh}) is usually traced, through Persian, to Arabic \emph{fatḥ}, \textquotedblleft{}an opening\textquotedblright{}, and so the taking of a city. It is feminine.
\end{cousinweb}

\begin{grammarlens}[title={What you've built}]
One more word for telling what happened.
\end{grammarlens}

\subsection*{Your turn}
\begin{itemize}
\raggedright
  \item \emph{Say it:} \emph{fateh}
  \item \emph{Say it:} \emph{fateh}, once more
  \item \emph{Say it:} \emph{fateh ke bād jashn huā.} — After the victory there was a celebration.
  \item \emph{Recall:} say the Urdu for to defeat, then the Urdu for to throw down, then say \emph{fateh} again
\end{itemize}

\begin{tcolorbox}[breakable,colback=teal!4,colframe=teal!35!black,title={Before you move on}]
What does \ur{فتح} mean? (\textquotedblleft{}A victory\textquotedblright{}.) Say it once more.
\end{tcolorbox}
\section[\texorpdfstring{\ur{شہید} (shahīd)}{شہید (shahīd)}]{\ur{شہید} (shahīd) — a martyr}

\label{lesson:UR-C311-shahid}



\begin{quote}
Before the new one: say the Urdu for to throw down, then the Urdu for a victory.
\end{quote}

\subsection*{You'll want to know: \ur{شہید}}
\textbf{\ur{شہید}} — \emph{shahīd} — \textquotedblleft{}a martyr\textquotedblright{}.

\begin{quote}
\ur{جنگ} \ur{میں} \ur{دو} \ur{سو} \ur{سپاہی} \ur{شہید} \ur{ہوئے۔}
\end{quote}

\begin{quote}
\emph{jang meṅ do sau sipāhī shahīd hue.}
\end{quote}

\begin{quote}
Two hundred soldiers fell in the war.
\end{quote}

Masculine or feminine, by the person meant.

\begin{cousinweb}
\textbf{\ur{شہید}} (\emph{shahīd}) is usually traced, through Persian, to Arabic \emph{shahīd}, \textquotedblleft{}a witness\textquotedblright{}, and so one who dies for a cause. \textbf{\ur{شہید} \ur{ہوئے}} is said of soldiers killed in battle.
\end{cousinweb}

\begin{grammarlens}[title={What you've built}]
One more word for telling what happened.
\end{grammarlens}

\subsection*{Your turn}
\begin{itemize}
\raggedright
  \item \emph{Say it:} \emph{shahīd}
  \item \emph{Say it:} \emph{shahīd}, once more
  \item \emph{Say it:} \emph{jang meṅ do sau sipāhī shahīd hue.} — Two hundred soldiers fell in the war.
  \item \emph{Recall:} say the Urdu for to throw down, then the Urdu for a victory, then say \emph{shahīd} again
\end{itemize}

\begin{tcolorbox}[breakable,colback=teal!4,colframe=teal!35!black,title={Before you move on}]
What does \ur{شہید} mean? (\textquotedblleft{}A martyr\textquotedblright{}.) Say it once more.
\end{tcolorbox}
\section[\texorpdfstring{\ur{صلح} (sulah)}{صلح (sulah)}]{\ur{صلح} (sulah) — a peace settlement; a reconciliation}

\label{lesson:UR-C311-sulah}



\begin{quote}
Before the new one: say the Urdu for a victory, then the Urdu for a martyr.
\end{quote}

\subsection*{You'll want to know: \ur{صلح}}
\textbf{\ur{صلح}} — \emph{sulah} — \textquotedblleft{}a peace settlement; a reconciliation\textquotedblright{}.

\begin{quote}
\ur{آخر} \ur{دونوں} \ur{بادشاہوں} \ur{میں} \ur{صلح} \ur{ہوئی۔}
\end{quote}

\begin{quote}
\emph{ākhir donoṅ bādshāhoṅ meṅ sulah huī.}
\end{quote}

\begin{quote}
In the end the two kings made peace.
\end{quote}

Feminine.

\begin{cousinweb}
is usually traced, through Persian, to Arabic \emph{ṣulḥ}, \textquotedblleft{}peace, a settlement\textquotedblright{}. \textbf{\ur{امن}} is peace as a state; \textbf{\ur{صلح}} is the peace two sides make.
\end{cousinweb}

\begin{grammarlens}[title={What you've built}]
One more word for telling what happened.
\end{grammarlens}

\subsection*{Your turn}
\begin{itemize}
\raggedright
  \item \emph{Say it:} \emph{sulah}
  \item \emph{Say it:} \emph{sulah}, once more
  \item \emph{Say it:} \emph{ākhir donoṅ bādshāhoṅ meṅ sulah huī.} — In the end the two kings made peace.
  \item \emph{Recall:} say the Urdu for a victory, then the Urdu for a martyr, then say \emph{sulah} again
\end{itemize}

\begin{tcolorbox}[breakable,colback=teal!4,colframe=teal!35!black,title={Before you move on}]
What does \ur{صلح} mean? (\textquotedblleft{}A peace settlement; a reconciliation\textquotedblright{}.) Say it once more.
\end{tcolorbox}
